\documentclass[10pt,a4paper,twocolumn]{article}

\usepackage[a4paper,top=1.75cm,bottom=2.0cm,left=1.75cm,right=1.75cm,columnsep=0.65cm]{geometry}
\usepackage{times}
\usepackage[T1]{fontenc}
\usepackage[utf8]{inputenc}
\usepackage{microtype}
\usepackage{amsmath}
\usepackage{booktabs}
\usepackage{tabularx}
\usepackage{enumitem}
\usepackage{xcolor}
\usepackage{natbib}
\usepackage[hidelinks]{hyperref}
\usepackage{titlesec}

\setlength{\parindent}{1em}
\setlength{\parskip}{0pt}
\setlength{\textfloatsep}{8pt plus 2pt minus 2pt}
\setlength{\floatsep}{7pt plus 2pt minus 2pt}
\setlength{\intextsep}{7pt plus 2pt minus 2pt}
\setlist{nosep,leftmargin=1.4em}
\titleformat{\section}{\large\bfseries}{\thesection}{0.55em}{}
\titleformat{\subsection}{\normalsize\bfseries}{\thesubsection}{0.45em}{}
\titlespacing*{\section}{0pt}{1.4ex plus .4ex}{0.7ex}
\titlespacing*{\subsection}{0pt}{1.0ex plus .3ex}{0.4ex}

\title{Reproducing LLMLingua: Evaluating Accuracy--Efficiency Trade-offs in Long-Context Prompt Compression}

\author{
CHEN IAT HONG \quad MC654102\\
KONG KIN MENG \quad MC652817\\
CHE CHON HONG \quad MC551841\\
\textit{University of Macau}\\
\textit{CISC7021 Applied Natural Language Processing -- Course Project Final Report}
}

\date{}

\begin{document}
\maketitle
\vspace{-0.8cm}

% ===========================================================================
% DRAFT STATUS: Abstract and Section 1 are written. Everywhere else the body
% is still the placeholder line
%     (text text text)
% with a "% TODO" note describing what must be written there.
% ===========================================================================

\begin{abstract}
Long prompts improve retrieval-augmented and in-context question answering, but
they also raise memory use, latency, and cost, and models may ignore information
placed in the middle of a long context. Prompt compression attacks both problems
without retraining the target model, yet its reported gains are usually measured
with proprietary LLMs and reported as token reductions rather than end-to-end
cost. We reproduce and evaluate four prompt settings on the full 150-example
LongBench \texttt{multifieldqa\_en} split with a 3B open-weight target model:
an uncompressed baseline, head-and-tail truncation, LLMLingua, and a
question-aware long-context variant of LLMLingua, the last two at nominal $2\times$
and $4\times$ compression. All settings share one fixed sample index, one prompt
template, and one Exact Match / token-F1 implementation. Accuracy decreases
monotonically with the compression ratio in every compressed family. Truncation
keeps most of the baseline quality (F1 0.4065 and 0.3620 against 0.4296), whereas
LLMLingua is markedly worse (0.2626 and 0.2207); replacing its chunked Phi-2
compressor with a single-pass question-aware compressor recovers part of the loss
(0.3070 and 0.2794) but does not change the ranking. In a same-machine rerun of
the first three settings, only truncation actually reduces end-to-end latency, by
44\% and 66\%, while LLMLingua at $2\times$ is 44\% slower than the uncompressed
baseline despite removing half of the prompt tokens.
\end{abstract}

% ---------------------------------------------------------------------------
\section{Introduction}

\subsection{Motivation}

Large language models (LLMs) are now routinely asked to answer a question from
several retrieved documents, to summarise a long article, or to learn a task from
many in-context demonstrations. Each of these settings pushes thousands of tokens
into a single prompt. The cost is not only monetary: a longer prompt occupies more
GPU memory, needs more prefill computation, and therefore increases the
time-to-first-token and the end-to-end latency of the serving system. Retrieved
context is also repetitive. Documents overlap with one another and with the
instruction, so a non-trivial fraction of the prompt carries little new
information for the question at hand. Finally, a large advertised context window
does not imply that the model uses it uniformly: \citet{liu-etal-2024-lost} show
that accuracy is typically highest when the relevant passage sits at the
beginning or the end of the context and drops when it sits in the middle.

Prompt compression addresses these issues by shortening the text that reaches the
target model while trying to preserve the information the task needs. Unlike
model compression or distillation, it requires neither retraining nor access to
the target model's weights, so the same compressor can be placed in front of any
open or proprietary endpoint. This makes it an attractive engineering option for
resource-constrained deployments, and it is the reason a family of
perplexity-based methods has been developed around it.

What is missing is not a new method but a clean measurement. The headline numbers
in this literature are typically obtained with proprietary target models and are
reported as reductions in input tokens. Token reduction, however, is only one side
of the ledger: the compressor itself is a language model that must run on the same
hardware, and its cost can exceed the generation time it saves. It is therefore
unclear whether the accuracy--efficiency trade-off reported for large proprietary
models still holds for a smaller open-weight model that a student project, or a
small deployment, can actually afford to run---and whether the compressor pays for
itself once its own latency is counted.

\subsection{Literature Review}

\citet{li-etal-2025-prompt-survey} organise the field and separate methods that
delete or rewrite visible prompt tokens from methods that encode the prompt into
learned continuous representations. The former are directly applicable to any
frozen model and are the family studied here.

Selective Context \citep{li-etal-2023-compressing} removes lexical units whose
self-information is low, showing that natural-language prompts contain
substantial redundancy. Its importance estimate is not conditioned on the
downstream question, so it may discard details that are predictable in general but
essential for a specific query.

LLMLingua \citep{jiang-etal-2023-llmlingua} introduces a coarse-to-fine framework
in which a small language model scores prompt components and individual tokens.
A budget controller first allocates different compression budgets to instructions,
demonstrations, and questions; iterative token-level compression then removes
low-value tokens in several passes, and distribution alignment reduces the gap
between the compressor and the target model. Reported gains reach roughly
$20\times$ compression with limited quality loss on some tasks.

LongLLMLingua \citep{jiang-etal-2024-longllmlingua} extends the idea to
long-context question answering, where useful evidence is sparse and
question-dependent. It adds question-aware document ranking, contrastive
perplexity, dynamic budget allocation, document reordering, and subsequence
recovery, and reports large token and cost savings on a multi-document Natural
Questions variant \citep{kwiatkowski-etal-2019-natural}.

Two evaluation resources complete the picture. \citet{liu-etal-2024-lost}
motivate the accuracy side of our measurement by documenting position bias, and
LongBench \citep{bai-etal-2024-longbench} supplies the long-context
question-answering data we use, in a format that requires no additional
preprocessing.

\subsection{Research Questions}

We ask three questions about a small, reproducible configuration:

\begin{itemize}
    \item \textbf{RQ1:} How does the compression ratio affect question-answering
    accuracy?
    \item \textbf{RQ2:} At the same input-token budget, does LLMLingua outperform
    simple truncation?
    \item \textbf{RQ3:} Does compression reduce end-to-end latency once the
    compression overhead is included?
\end{itemize}

\subsection{Contributions}

\begin{itemize}
    \item A shared evaluation harness, used unchanged by all four experiments: one
    fixed index of the 150 official \texttt{multifieldqa\_en} test examples, one
    chat-template prompt renderer, one set of Exact Match / token-F1
    implementations, and one configuration module.
    \item Four compression settings implemented end to end and run to completion on
    the full split: uncompressed, head-and-tail truncation, LLMLingua, and a
    question-aware long-context variant of LLMLingua, the latter two at nominal
    $2\times$ and $4\times$ compression (seven result settings in total).
    \item A paired, per-sample comparison with 95\% confidence intervals, which is
    more informative than comparing leaderboard averages and lets us state
    explicitly which differences survive noise.
    \item A single-machine rerun of the five uncompressed/truncation/LLMLingua
    settings, so that the latency question (RQ3) is answered on one GPU rather
    than by comparing four different machines.
    \item A negative result with a diagnosed cause: LLMLingua loses to plain
    truncation at an equal token budget, chunking explains only part of that gap,
    and the residual gap persists after the chunking defect is removed.
\end{itemize}

% ---------------------------------------------------------------------------
\section{Detailed Description}
(text text text)

\subsection{Task and Dataset}
(text text text)

% TODO: LongBench multifieldqa_en; question + retrieved documents + gold answers;
% why all 150 official test examples were kept; how the fixed index is stored in
% shared/samples.json (seed 42).

\subsection{Models and Hardware}
(text text text)

% TODO: target model Llama-3.2-3B-Instruct (gated id vs unsloth mirror); 128k
% native context and why a 32k prompt budget was self-imposed; compressor models
% (microsoft/phi-2 in Part 3; Qwen2.5-1.5B-Instruct in Part 4); the four machines
% and the single-machine rerun used for latency.

\subsection{Shared Evaluation Harness}
(text text text)

% TODO: one contract for all members — files, what is frozen, what a member may
% change. Expand on the four sub-blocks below.

\subsubsection{Sample Index}
(text text text)

% TODO: 150 ids, order, seed, why cross-part comparability depends on it.

\subsubsection{Prompt Template}
(text text text)

% TODO: the chat template and instruction/context formatting; why prompt_tokens
% counts the chat-templated prompt; the token-budget fallback.

\subsubsection{Metrics}
(text text text)

% TODO: SQuAD-style Exact Match and token-level F1; mean over 150 samples; what
% EM/F1 ignore (e.g. LaTeX fragments, truncation at 64 new tokens).

\subsubsection{Configuration}
(text text text)

% TODO: greedy decoding, max_new_tokens = 64, MAX_PROMPT_TOKENS, dtype, seeds.

\subsection{Compared Methods}
(text text text)

% TODO: one paragraph per method plus the target ratio definitions (document
% ratio vs whole-prompt ratio).

\subsubsection{Part 1: Uncompressed Baseline}
(text text text)

% TODO: mechanism, parameters, expected role as reference.

\subsubsection{Part 2: Head-and-Tail Truncation}
(text text text)

% TODO: budget b = round(n / r); head_frac = 0.5; instruction and question are
% never truncated; why this is a non-semantic baseline.

\subsubsection{Part 3: LLMLingua}
(text text text)

% TODO: coarse-to-fine, iterative token filtering, rates 0.5 / 0.25, phi-2 as
% compressor, measurement tokenizer.

\subsubsection{Part 4: LLMLingua-Enhanced}
(text text text)

% TODO: what changed relative to Part 3 (single-pass Qwen2.5-1.5B-Instruct,
% question-aware ranking, sentence-level filter), and why the improvement cannot
% be attributed to a single factor.

\subsection{Implementation}
(text text text)

% TODO: how the pipeline is structured end to end; the three sub-blocks below.

\subsubsection{Two-Stage Compress and Evaluate Pipeline}
(text text text)

% TODO: compressor and target model never co-resident; two environments because
% of the transformers version conflict; JSONL checkpoints and resume.

\subsubsection{Chunked Compression and Chunked Prefill}
(text text text)

% TODO: 1,400-token compressor chunks in Part 3 (296/300 rows chunked); no
% chunking in Part 4; --chunk-size / --chunk-threshold prefill in Part 2; how
% chunking is recorded per row.

\subsubsection{Memory and Out-of-Memory Handling}
(text text text)

% TODO: Colab T4 limits, OOM retry with halved prompt budget, skipped-and-retried
% rows, oom_skipped_ids.

\subsubsection{System Overview}
(text text text)

% TODO: insert the pipeline / component diagram here (figure with caption and
% label). One figure covering data -> (compression) -> generation -> scoring.

\subsection{Experimental Setup}
(text text text)

% TODO: which settings were run, on which machine, at which ratio, when, and
% which result file each row comes from.

\subsection{Evaluation Protocol}
(text text text)

% TODO: primary metrics; paired-by-id F1 differences; 95% CI = mean +/- 1.96 SE;
% how "significant" is decided; which contrasts are reported; why per-part
% latencies are not compared across machines.

% ---------------------------------------------------------------------------
\section{Results and Evaluation}
(text text text)

\subsection{Main Results}
(text text text)

% TODO: the seven-setting table (uncompressed, head-tail 2x/4x, LLMLingua
% 2x/4x, LLMLingua-enhanced 2x/4x): n, EM, F1, prompt tokens, token reduction,
% compression / generation / total latency. Reference part4 results_table.md.

\subsection{RQ1: Accuracy versus Compression Ratio}
(text text text)

% TODO: monotone F1 decline from 1x to 2x to 4x in all three compressed families;
% include the accuracy--compression figure (figures/accuracy_compression_tradeoff).

\subsection{RQ2: LLMLingua versus Truncation at Equal Token Budget}
(text text text)

% TODO: paired contrasts — LLMLingua loses to truncation (-0.1440 / -0.1414);
% enhanced LLMLingua improves over LLMLingua (+0.0444 / +0.0587) but still loses
% to truncation (-0.0995 / -0.0827); all CIs exclude zero.

\subsection{RQ3: End-to-End Latency}
(text text text)

% TODO: use the single-machine rerun only; truncation cuts total latency by
% 43.6% / 65.7%; LLMLingua 2x is 43.9% slower than uncompressed once compression
% is charged; Part 4 has no common-hardware rerun (state this explicitly).

\subsection{Qualitative Analysis}
(text text text)

% TODO: reproducible success and failure cases from
% part3_llmlingua/qualitative_cases.md; what compression removes that the
% question needed; what survives intact.

\subsection{Discussion}
(text text text)

% TODO: why perplexity-based token removal underperforms a dumb baseline here
% (fragmentation, names/numbers, 3B target model, short prompts, question not
% used for budget allocation); what this implies about the paper's claims.

\subsection{Threats to Validity}
(text text text)

% TODO: per-part hardware differences; stale results_same_machine export;
% differing verdict for Uncompressed -> Head-tail 2x between runs; single
% dataset, single model, single head_frac; Part 4 not in the rerun.

% ---------------------------------------------------------------------------
\section{Conclusion}
(text text text)

\subsection{Summary of Findings}
(text text text)

% TODO: answer RQ1--RQ3 in three sentences, each with its number.

\subsection{Lessons Learned}
(text text text)

% TODO: about reproduction, about evaluation contracts, about measuring
% end-to-end cost rather than token counts.

\subsection{Problems and Difficulties}
(text text text)

% TODO: gated model access; Windows console encoding crash at sample 87; the
% compressor context limit forcing chunking; Colab T4 OOM forcing chunked
% prefill; the Windows commit-limit failure that blocked Part 4 locally.

\subsection{Limitations}
(text text text)

% TODO: scope and measurement limits, one bullet per limitation.

\subsection{Future Work}
(text text text)

% TODO: answer-position failure analysis; LongLLMLingua-style reordering;
% larger target models; more ratios and datasets; latency measurement on one
% common GPU for all four parts.

% ---------------------------------------------------------------------------
\section{Use of Generative AI}
(text text text)

% TODO: which tools were used for what (literature organisation, drafting,
% formatting, code assistance), and how numbers and claims were verified.

% ---------------------------------------------------------------------------
% References: carried over from report/acl_latex.tex. Keep only the entries that
% are actually cited, and add any new ones in the same format.
\begin{thebibliography}{9}

\bibitem[Bai et al.(2024)]{bai-etal-2024-longbench}
Yushi Bai, Xin Lv, Jiajie Zhang, Hongchang Lyu, Jiankai Tang, Zhidian Huang, Zhengxiao Du, Xiao Liu, Aohan Zeng, Lei Hou, Yuxiao Dong, Jie Tang, and Juanzi Li. 2024.
\newblock LongBench: A bilingual, multitask benchmark for long context understanding.
\newblock In \emph{Proceedings of the 62nd Annual Meeting of the Association for Computational Linguistics (Volume 1: Long Papers)}, pages 3119--3137.

\bibitem[Jiang et al.(2023)]{jiang-etal-2023-llmlingua}
Huiqiang Jiang, Qianhui Wu, Chin-Yew Lin, Yuqing Yang, and Lili Qiu. 2023.
\newblock LLMLingua: Compressing prompts for accelerated inference of large language models.
\newblock In \emph{Proceedings of the 2023 Conference on Empirical Methods in Natural Language Processing}, pages 13358--13376.

\bibitem[Jiang et al.(2024)]{jiang-etal-2024-longllmlingua}
Huiqiang Jiang, Qianhui Wu, Xufang Luo, Dongsheng Li, Chin-Yew Lin, Yuqing Yang, and Lili Qiu. 2024.
\newblock LongLLMLingua: Accelerating and enhancing LLMs in long context scenarios via prompt compression.
\newblock In \emph{Proceedings of the 62nd Annual Meeting of the Association for Computational Linguistics (Volume 1: Long Papers)}, pages 1658--1677.

\bibitem[Kwiatkowski et al.(2019)]{kwiatkowski-etal-2019-natural}
Tom Kwiatkowski, Jennimaria Palomaki, Olivia Redfield, Michael Collins, Ankur Parikh, Chris Alberti, Danielle Epstein, Illia Polosukhin, Jacob Devlin, Kenton Lee, Kristina Toutanova, Llion Jones, Matthew Kelcey, Ming-Wei Chang, Andrew M. Dai, Jakob Uszkoreit, Quoc Le, and Slav Petrov. 2019.
\newblock Natural Questions: A benchmark for question answering research.
\newblock \emph{Transactions of the Association for Computational Linguistics}, 7:453--466.

\bibitem[Li et al.(2023)]{li-etal-2023-compressing}
Yucheng Li, Bo Dong, Frank Guerin, and Chenghua Lin. 2023.
\newblock Compressing context to enhance inference efficiency of large language models.
\newblock In \emph{Proceedings of the 2023 Conference on Empirical Methods in Natural Language Processing}, pages 6342--6353.

\bibitem[Li et al.(2025)]{li-etal-2025-prompt-survey}
Zongqian Li, Yinhong Liu, Yixuan Su, and Nigel Collier. 2025.
\newblock Prompt compression for large language models: A survey.
\newblock In \emph{Proceedings of the 2025 Conference of the Nations of the Americas Chapter of the Association for Computational Linguistics: Human Language Technologies (Volume 1: Long Papers)}, pages 5807--5823.

\bibitem[Liu et al.(2024)]{liu-etal-2024-lost}
Nelson F. Liu, Kevin Lin, John Hewitt, Ashwin Paranjape, Michele Bevilacqua, Fabio Petroni, and Percy Liang. 2024.
\newblock Lost in the middle: How language models use long contexts.
\newblock \emph{Transactions of the Association for Computational Linguistics}, 12:157--173.

\end{thebibliography}

\end{document}
